\chapter{A Short History of Agents, Control, and Commitment}
\label{ch:history}

In 1979, a database researcher worried about a bank transfer that debits one
account and crashes before crediting the other. In 2026, an engineer worries about
an agent that charges a card and crashes before recording the order. These are the
same problem, and only one of those people had a protocol for it.

Modern agents look new. They are language-driven, tool-using, autonomous, adaptive,
and they arrived fast enough that the field has mostly been reinventing rather than
reading. But their characteristic failures were catalogued and largely solved in other
disciplines, some of them decades ago. Runaway resource consumption. Oscillation.
Acting on stale state. Committing to work that cannot be undone.

This chapter is a tour of that inheritance. It is deliberately incomplete; the
point is to show where each of the next twelve chapters got its ideas, and to make
the case that the failures you are seeing in production are older than they look.

It is also the one chapter you can skip. Nothing later depends on it, and each
borrowed mechanism is developed properly at the point of use. Terms are glossed
where they first appear, so no prior background in control theory or transaction
processing is assumed.
\section{Part I: Computation Without Action}

Classical computer science is built on a precise abstraction. Computation as a
closed process that transforms input into output and then halts. This abstraction
has been extraordinarily successful. It underlies algorithm design, complexity
theory, compilers, and formal verification. It also excludes, by construction,
nearly every phenomenon that dominates the behavior of modern agent systems.

The abstraction is sound. The trouble is that it gets applied outside its domain
without anyone noticing the boundary was crossed.

In the Turing model, correctness is defined relative to termination and final
output \cite{turing1936computable}. A computation that does not halt has no semantic
value, and a computation that halts with the wrong output is incorrect. Partial
progress, intermediate states, and resource consumption during execution have no
formal status. This notion of correctness is appropriate for algorithms. It is
incompatible with agents. An agent that halts is typically failing, while an agent
that continues operating under constraints may be succeeding. For agents,
correctness must be defined incrementally, over time, rather than at termination.

Equally important is the absence of action. A Turing machine modifies symbols on a
tape. These modifications have no effect beyond the computation itself. Restarting
the computation erases all consequences. As a result, classical computation has no
need for an explicit notion of commitment. There is no distinction between
tentative and final state, because all state is internal and disposable. This
assumption is harmless for algorithms. It becomes dangerous when systems interact
with the external world, where actions may be irreversible and observable by other
systems.

Cost in classical computation is comparative and asymptotic. Time and space
complexity order algorithms by growth rate; they do not constrain execution.
Exceeding an expected runtime bound does not invalidate a computation. It merely
makes it inefficient. Agent systems operate under absolute limits, namely deadlines,
quotas, budgets, and escalation thresholds. Exceeding these limits terminates
execution or triggers control paths that cannot be ignored. Classical cost models
do not express this distinction. When they are applied to agents, budget exhaustion
is treated as a performance issue rather than as a correctness failure.

Because classical computation lacks absolute cost limits, it also lacks budget
semantics. There is no notion of refusing to compute because resources are
insufficient, nor of reserving resources for future steps. Every computation is
implicitly authorized to consume unbounded resources until termination. When this
assumption is inherited by agent systems, unbounded verification, retrieval, and
retry loops become the default behavior rather than a design error.

Another implicit assumption is reversibility. Classical computation assumes that
errors can be corrected by restarting execution. State corruption does not
propagate beyond the failed computation. Agents do not enjoy this property. They
accumulate state across decisions, interact with shared environments, and
influence other systems. Errors propagate through memory, configuration, and
downstream actions. Restarting does not restore the initial conditions. Any design
that relies on implicit rollback is therefore unsound.

Notably, none of these assumptions were fundamentally altered by the introduction
of interactivity. Time-sharing systems added preemption and scheduling
\cite{corbato1962experimental}, but they did not change the semantics of execution.
Processes could be delayed or terminated, yet their actions remained largely
reversible and local. Interactivity introduced latency constraints without
introducing commitment semantics. The system still did not act in an irreversible
environment.

Despite the evolution of computing toward networked services and automation, these
assumptions persist. Many modern systems still implicitly assume that execution is
cheap relative to decision-making, that errors can be retried safely, and that
state can be discarded without consequence. These assumptions are reasonable for
algorithms. They are false for autonomous agents with authority.

Agent systems violate every foundational assumption of classical computation. They
do not halt cleanly, they act irreversibly, they operate under hard budgets, and
they accumulate state across decisions. Treating agents as extended algorithms
therefore produces systems that are formally well-defined yet operationally
unsafe. The remainder of this book introduces the semantics that classical computation deliberately omits. Budgets, gates, commitment, abstention, and refusal all compensate for that gap.

\section{Part II: Control Systems Before Intelligence}

The first systems that resemble modern agents did not reason symbolically, learn
from data, or optimize explicit objectives. They regulated physical processes.
Their primary concern was not intelligence but stability. This historical fact is
easy to overlook, yet it explains why many contemporary agent failures look less
like reasoning errors and more like control failures.

Control systems were designed to operate continuously, under uncertainty, and in
the presence of noise. They sensed state, acted on the environment, and observed
the consequences of those actions. Crucially, they did so without assuming perfect
models or complete information. These properties place control systems much closer
to deployed agents than to classical algorithms.

The canonical example is Watt’s centrifugal governor, which regulated steam engine
speed via mechanical feedback. The governor continuously adjusted fuel input based
on observed rotational velocity, stabilizing the system without computing an
optimal solution. The design objective was not maximum efficiency but bounded
behavior \cite{bennett1979history}. Overshoot and oscillation were treated as
failures even if average performance was high.

This emphasis on stability introduced a key distinction. Optimality is secondary
to safety. A controller that produces slightly suboptimal output but remains
stable is preferable to one that achieves higher average performance while
occasionally destabilizing the system. This trade-off reappears in agent design as
abstention, refusal, and conservative gating. Systems that always act aggressively
fail less often in simulations but more often in production.

Control theory formalized these ideas through feedback loops. A controller observes
the current state, compares it to a desired operating region, and applies a
corrective action. Importantly, the controller does not assume that its actions
will have the intended effect. It continuously corrects based on observed
outcomes. This contrasts sharply with planning-centric agent designs that assume
execution follows intention unless proven otherwise.

Another critical contribution of control theory is the treatment of delay and
latency as first-class concerns. Delayed feedback can destabilize an otherwise
sound controller. As a result, control systems are often deliberately conservative,
trading responsiveness for stability. This logic directly contradicts the
impulse to speculate aggressively in agent systems without accounting for commit
latency or verification delay.

Cybernetics generalized these ideas beyond mechanical systems. Wiener emphasized
that intelligent behavior emerges from feedback, not from complex internal
representation alone \cite{wiener1948cybernetics}. Systems must be evaluated by
their trajectories, not by isolated outputs. A system that occasionally produces
excellent responses but diverges over time is inferior to one that behaves
predictably within bounds.

Notably, early control systems treated failure as inevitable. Noise, drift, and
disturbance were expected. Controllers were therefore designed to degrade
gracefully rather than to achieve perfect performance. This assumption is absent
from many modern agent designs, which implicitly assume that reasoning failures are
rare and that recovery is always possible.

The contrast between control systems and symbolic AI is instructive. Control
systems act continuously, assume uncertainty, and prioritize stability. Symbolic
AI systems plan discretely, assume correct execution, and optimize objectives.
Modern agents inherit characteristics of both, but failures arise when the latter
assumptions dominate.

When agents are designed without explicit feedback control, they behave like open-
loop systems. Actions are chosen based on internal reasoning and executed with
minimal correction. In environments with irreversible effects, such designs are
fragile. Small errors accumulate until recovery is impossible.

The control-theoretic perspective therefore introduces several primitives that
reappear throughout this book, such as explicit feedback, bounded response, conservative
action, and stability over optimality. These primitives are not add-ons. They are
preconditions for safe operation in uncertain environments.

Nobody is suggesting you build your agent as a PID controller. That is the standard industrial feedback controller, which corrects using the present error, its accumulated history, and its rate of change. The transferable
lesson is about ordering. Anything that acts repeatedly on the world has to be
designed as a control system first and an intelligent one second. Reverse that
order and you get what the field keeps producing, systems that demo beautifully
and misbehave the moment they are embedded in a real process with real feedback
delays.

\section{Part III: Planning Systems and the Execution Gap}

Classical AI approached agency through planning. An agent was modeled as a system that constructs a plan, meaning a sequence of actions leading from an initial state to a goal state, and then executes it. This formulation made reasoning explicit and
symbolic, but it embedded assumptions about execution that fail in real
environments.

Early planning systems such as STRIPS described actions by their preconditions and effects. They assumed deterministic actions, complete state observability, and negligible execution cost \cite{fikes1971strips}. A plan
was correct if, when executed exactly as specified, it reached the goal state.
Execution was treated as a mechanical process whose only responsibility was to
follow the plan.

This separation between planning and execution created a structural blind spot.
Plans were optimized for logical correctness, not for robustness under deviation.
Unexpected events invalidated plans, requiring replanning from scratch. In
environments where deviation is common, this leads to continual replanning and
rapid cost escalation.

The assumption that execution is cheap is particularly damaging. Planning systems
implicitly authorize arbitrary sequences of actions without regard to cost,
latency, or side effects. When such systems are embedded in real environments,
verification, retries, and recovery actions accumulate along execution traces,
often dominating total cost.

The execution gap widens further when plans are long-horizon. Errors early in a
plan propagate forward, constraining later options. Because planning systems
evaluate correctness at the end of execution, they provide no mechanism for early
intervention when a plan is becoming unsafe or unaffordable.

Attempts to repair this gap often introduced monitoring and replanning. The agent
would observe discrepancies between expected and actual state and trigger a new
planning phase. While this improves logical correctness, it does not address the
underlying cost structure. Replanning itself consumes resources, and repeated
monitor–replan cycles can place the system on an expensive critical path.

A deeper issue is that planning systems conflate intention with authorization.
Producing a plan implicitly authorizes its actions. There is no intermediate
decision point at which actions can be refused, deferred, or bounded based on
policy, cost, or risk. Once a plan is accepted, execution proceeds unless a failure
is detected.

This model is incompatible with environments where actions are irreversible or
policy-constrained. An action that violates a safety constraint cannot be undone
by replanning. An action that exhausts a budget cannot be recovered by generating a
new plan.

The reliance on idealized world models exacerbates these problems. Planning systems
assume that all relevant preconditions and effects are known and encoded. In real
systems, preconditions are often probabilistic, hidden, or adversarial. Plans
constructed under incomplete models appear valid but fail at execution time in
systematic ways.

The planning literature gradually recognized these limitations, introducing
contingent planning, execution monitoring, and replanning under uncertainty.
However, these extensions retained the core assumption that planning is the
primary locus of control. Execution remained subordinate.

Modern agent systems often repeat this pattern. Large language models generate
plans or chains of reasoning, which are then executed by tools. Failures are
attributed to poor reasoning rather than to the absence of explicit execution
control. Verification and safety checks are bolted on after the fact, leading to
cascades and brittle behavior.

Blaming the planner here is a diagnostic error. The execution gap opens because
execution was treated as a detail to be handled later, and it closes only when the
system makes control decisions at execution time, including act, verify, abstain, or refuse.
Those four decisions depend on what the world looks like at the moment of acting,
which is information a plan built in advance does not contain. Chapter~\ref{ch:gates}
is entirely about deriving them.

This book adopts the opposite ordering. Plans and proposals are treated as
tentative. Authorization, cost, and policy are enforced at execution time, not
assumed at planning time. The mechanisms introduced in later chapters, gates,
budgets, and commit protocols, exist to close the execution gap that planning-based
agent models leave open.

\section{Part IV: Learning Systems and Soft Constraints}

Reinforcement learning reframed agents as systems that learn by interacting with an
environment, selecting actions to maximize cumulative reward
\cite{sutton2018reinforcement}. This formulation addressed a core limitation of
classical planning. It treated action and consequence as inseparable. However, it
introduced a different abstraction gap, one that becomes critical in deployed
agent systems.

The defining move of reinforcement learning is the reduction of all objectives to
a scalar reward. Performance, safety, cost, and policy compliance are encoded, if
at all, as components of a single objective function. This reduction is convenient
for optimization, but it collapses distinctions that matter operationally. In
particular, it replaces hard constraints with soft trade-offs.

In a reward-based formulation, violating a safety constraint is not categorically
forbidden; it is merely assigned a large negative reward. Whether the agent
violates the constraint becomes a matter of optimization rather than admissibility.
This is acceptable in simulated environments where violations can be sampled,
averaged, and penalized. It is unacceptable in environments where violations are
irreversible or catastrophic.

A related consequence is the averaging away of rare failures. Reinforcement
learning optimizes expected return. Low-probability, high-impact failures
contribute little to the expectation unless explicitly amplified. As a result,
policies that are optimal in expectation may still produce unacceptable tail
behavior. This phenomenon is well documented in safety-critical domains, where
systems perform well until a rare event triggers disproportionate damage
\cite{amodei2016concrete}.

Attempts to address these issues introduced constrained Markov decision processes,
which separate rewards from constraints \cite{altman1999constrained}. While this
formalism restores the distinction between objectives and constraints, it does not
eliminate the reliance on expectation. Constraints are typically enforced in
expectation or on average over trajectories, not at every step.

For agents that act in the world, step-wise constraint enforcement is often
required. An action that violates a policy cannot be justified by compensating
future rewards. The action must simply not occur. This requirement lies outside
the expressive scope of standard reinforcement learning formulations.

Another limitation arises from the treatment of cost. In many learning systems,
cost is incorporated as a negative reward. This encourages cost-efficient behavior
on average but provides no mechanism for enforcing hard budgets. An agent may
learn to respect cost constraints most of the time while still occasionally
exhausting its budget. In operational systems, such occasional violations are
unacceptable.

Learning systems also assume that policy improvement is the primary path to
correctness. Failures are treated as data points from which the agent can learn.
This assumption presumes that failures are survivable and that the system can
continue operating after them. Agents with real authority do not enjoy this
luxury. Some failures terminate the system or trigger irreversible consequences,
leaving no opportunity for learning.

Modern agent architectures often combine learned policies with heuristic guards
and filters. While this improves safety in practice, it obscures responsibility.
When a system fails, it is unclear whether the failure arose from the learned
policy, the reward specification, or the surrounding control logic. This
entanglement complicates analysis and certification.

The problem is where the learning sits. Learning is excellent at improving proposals. It suggests which actions might work, which evidence looks relevant, and which plans are plausible. It is a poor sole mechanism for enforcing an invariant, a budget, or a policy. Proposals are allowed to be wrong sometimes. Enforcement is
not.

This book therefore treats learning as a proposal mechanism rather than as a
control mechanism. Learned components generate candidates. Control logic, gates,
budgets, and commit protocols, determines what is executed. This separation
restores hard constraints without discarding the benefits of learning.

The historical lesson is that reward maximization alone is an insufficient
foundation for agency in real environments. Systems that rely on soft constraints
inevitably discover the constraints they violated only after the violation has
occurred. Agents that must operate safely cannot afford this mode of discovery.

\section{Part V: Commitment Is a Systems Problem}

The need for explicit commitment did not arise in artificial intelligence. It
arose in distributed systems, where computation interacts with shared state,
partial failure, and irreversible effects. The solutions developed in this domain
address problems that modern agent systems now encounter in different form.

Early distributed systems exposed a fundamental mismatch between computation and
action. Operations that appeared local from the perspective of a program had
global effects on shared state. Partial failure made it impossible to assume that
an operation either completed fully or not at all. This invalidated the implicit
reversibility assumed by classical computation.

Transaction processing systems responded by separating execution from commitment.
Two-phase commit introduced a protocol in which participants first prepare to
commit and only later make changes durable \cite{gray1978notes}. This separation
allowed systems to reason about tentative actions without exposing irreversible
effects prematurely. Crucially, commitment became an explicit, coordinated step
rather than an implicit consequence of execution.

What the transaction-processing community worked out, at some expense, is that
commitment is a systems-level property. It depends on coordination among
participants, on durability, and on a defined story for what happens when a
participant disappears mid-protocol. Each participant can be locally correct and
the system can still commit at the wrong moment and lose money.

That result transfers directly to agents invoking tools, editing configurations, or
kicking off external processes, and Chapter~\ref{ch:speculation} spends its length
on it.

Distributed systems also normalized failure. Nodes crash, messages are lost, and
timeouts occur routinely. As a result, correctness conditions were reformulated to
account for partial progress and aborted execution. Systems were designed to make
failure visible and manageable rather than exceptional.

Admission control provides a complementary lesson. Systems learned to reject work
early when resources were insufficient, rather than accepting work and failing
later. Refusal became a stability mechanism. A system that says “no” preserves
overall correctness better than one that says “yes” and collapses under load.

This runs against the instinct that a system should accept whatever work it can.
Operators learned otherwise the hard way. Unbounded acceptance produces cascading
failure, where the overloaded component's slowness becomes its upstream's queue
becomes an outage. Agents without an explicit refusal path inherit the whole
pattern, and Chapter~\ref{ch:gates} promotes refusal to a first-class outcome for
exactly this reason.

Congestion control makes the same point in a different register. TCP reads packet
loss as an instruction to slow down \cite{jacobson1988congestion}. A sender that
responds to failure by trying harder degrades the network for everyone, itself
included; a sender that backs off gets more throughput over any interval that
matters. Restraint is the throughput strategy. An agent that retries a failing tool
faster and faster is rediscovering congestion collapse on someone else's API.

Security systems further refined the notion of authorized action. Capability-based
security replaced ambient authority with explicit rights
\cite{saltzer1975protection}. Possession of a capability grants permission to act;
absence forbids it. Intent, reasoning, or good faith are irrelevant.

This model aligns closely with agent tool use. An agent should not execute an
action because it believes the action is reasonable. It should execute an action
only if it possesses the explicit capability to do so. Authorization precedes
reasoning.

Look at what these systems have in common. Execution is tentative. Commitment is
explicit. Refusal is routine. Failure is expected and planned for. None of them
optimize the best case; all of them bound the worst case, which is the trade you
make when other people depend on you.

Agent systems keep rediscovering these ideas one at a time and under different
names. A verification step is a prepare phase. A tool wrapper is a capability
check. A budget cap is admission control. A retry limit is backoff. Discovered
piecemeal and wired in ad hoc, they interact badly and fail in combination.
Specified explicitly, they compose.

Commitment, in particular, cannot be inferred from intent. It has to be earned
through protocol, and an agent that acts without commitment semantics has quietly assumed the Turing model in an environment where it does not hold. That model presumes disposable state, restartable execution, and no external consequences.

The last year has supplied an unusually literal demonstration. Agents built on
speculate-then-abort designs turned out to leave the aborted branch resident in the
serving session's key/value cache, where the model continues to attend to it, so a
branch the application believes it discarded goes on shaping behavior
\cite{abortedkv2026}. Checkpoint-and-rollback machinery, added for exactly this
kind of safety, has itself been shown to be attackable \cite{saferesume2026}. Both
findings say the same thing the transaction community said in the 1970s. An abort
that is not complete is not an abort, and completeness is a property you engineer
rather than assume.

This book takes the systems view throughout. Actions are proposed, evaluated, and
then committed through mechanisms you can point at. Refusal and abstention are
outcomes with their own thresholds. Commitment waits until the system can justify
it. None of this is conservatism for its own sake. It is what correctness costs
once actions are irreversible and failure is ordinary.

% ============================================================
% Bibliography: Historical Foundations of Control and Commitment
\section{Fallacies and Pitfalls}

\textbf{Fallacy. This is all just history.} Every mechanism in this chapter is
re-derived, badly, by teams that did not know it existed. The cost of not reading
it is measured in incidents, not in ignorance.

\textbf{Fallacy. Agents are a new kind of system.} They are a new kind of
\emph{component} inside a very old kind of system, one that acts on shared
state, under uncertainty, with irreversible effects.

\textbf{Fallacy. The classical model applies because the code is classical.} The
Turing model assumes disposable state, restartable execution, and no external
consequence. An agent violates all three, and inherits none of the guarantees.

\textbf{Fallacy. Refusing work is a degraded mode.} Admission control exists
because unbounded acceptance produces cascading failure. Saying no is how a
system stays available.

\textbf{Pitfall. Borrowing a mechanism without its assumptions.} Two-phase commit
assumes participants that can prepare and hold. Optimistic concurrency assumes
rare conflicts. A mechanism transplanted past its assumptions fails silently and
looks like the agent's fault.

\textbf{Pitfall. Reading control theory as a metaphor.} Feedback, bounded
response, and stability margins are quantitative properties. Used as vocabulary
rather than as design constraints, they explain failures without preventing them.

\section{Summary}

Five traditions, and what each one hands you.

Classical computation gave us correctness at termination, which is the wrong
frame for a system that is failing when it halts. It also gave us cost as an
asymptotic comparison rather than an absolute limit, which is why budget
exhaustion still gets filed as a performance issue instead of a correctness
failure.

Control theory contributed feedback, bounded response, and the preference for
stability over optimality, and the ordering claim that anything acting
repeatedly on the world is a control system before it is an intelligent one.

Planning research found the execution gap, namely a plan built in advance cannot contain
the information needed to decide, at the moment of acting, whether to act at all.

Learning research showed where learning belongs. It is excellent at proposing and
unsuitable as the sole enforcer of an invariant, because proposals may be wrong
and enforcement may not.

Distributed systems and transaction processing supplied the rest, such as execution as
tentative, commitment as an explicit protocol, refusal as a stability mechanism,
backoff as the throughput strategy, and capability as the basis of authorization.

The through-line is that every one of these disciplines converged on bounding the
worst case rather than optimizing the best one, and did so after paying for the
alternative.

\section*{Exercises}

\begin{enumerate}
\item Take a failure you have seen in an agent system and locate it in one of the
five traditions above. Name the mechanism that discipline would have applied.

\item TCP treats packet loss as a signal to slow down. Describe an agent retry
policy built on the same principle, and state what signal plays the role of
packet loss.

\item Admission control rejects work when resources are insufficient rather than
accepting and failing later. Write the admission test for an agent that must
reserve budget for a mandatory finishing step.

\item The execution gap says control decisions cannot be derived from a plan
alone. Give a concrete task where a correct plan still requires a runtime
decision to act, verify, abstain, or refuse, and say which of the four applies.

\item Two-phase commit separates prepare from commit. Map its two phases onto an
agent that proposes a database migration, and identify what plays the role of a
participant that disappears mid-protocol.
\end{enumerate}

\section*{Bibliographic Notes}

The termination-based notion of correctness originates with Turing
\cite{turing1936computable}; its mismatch with continuously operating systems is
the starting point for the reactive-systems literature.

Control theory's contribution here is Wiener's cybernetics \cite{wiener1948} and
the feedback tradition that followed; the stability-over-optimality preference is
standard in any control text and is the part most often lost in translation to
software.

Classical planning is represented by STRIPS \cite{fikes1971strips}; the execution
gap it exposes was addressed by later work on reactive and contingent planning.
Sutton and Barto \cite{sutton2018reinforcement} cover the learning tradition,
including why penalty terms are a poor substitute for hard constraints.

The transaction-processing lineage runs through Gray's work on commit protocols
\cite{gray1978notes} and the broader treatment in \cite{gray1992transaction};
sagas for long-running transactions are due to Garcia-Molina and Salem
\cite{garciamolina1987sagas}. Congestion control as restraint-for-throughput is
Jacobson \cite{jacobson1988congestion}. Capability-based authorization derives
from Saltzer and Schroeder \cite{saltzer1975protection}.

The two contemporary results cited in the closing section, incomplete abort
through retained cache state \cite{abortedkv2026} and the attackability of
checkpoint/rollback machinery \cite{saferesume2026}, are the clearest recent
demonstrations that these old lessons are still being relearned.

% ============================================================

\section*{Bibliography}

\begin{thebibliography}{20}

\bibitem{turing1936computable}
A.~M. Turing.
\newblock On computable numbers, with an application to the Entscheidungsproblem.
\newblock \emph{Proceedings of the London Mathematical Society}, 42(2):230--265,
1936.

\bibitem{knuth1997art}
D.~E. Knuth.
\newblock \emph{The Art of Computer Programming, Volume 1: Fundamental Algorithms}.
\newblock Addison--Wesley, 3rd edition, 1997.

\bibitem{corbato1962experimental}
F.~J. Corbató, M.~M. Daggett, and R.~C. Daley.
\newblock An experimental time-sharing system.
\newblock \emph{Proceedings of the AFIPS Fall Joint Computer Conference}, 1962.

\bibitem{bennett1979history}
S.~Bennett.
\newblock \emph{A History of Control Engineering, 1800--1930}.
\newblock IEE Control Engineering Series, 1979.

\bibitem{wiener1948cybernetics}
N.~Wiener.
\newblock \emph{Cybernetics. Or Control and Communication in the Animal and the
Machine}.
\newblock MIT Press, 1948.

\bibitem{fikes1971strips}
R.~E. Fikes and N.~J. Nilsson.
\newblock STRIPS: A new approach to the application of theorem proving to problem
solving.
\newblock \emph{Artificial Intelligence}, 2(3--4):189--208, 1971.

\bibitem{sutton2018reinforcement}
R.~S. Sutton and A.~G. Barto.
\newblock \emph{Reinforcement Learning. An Introduction}.
\newblock MIT Press, 2nd edition, 2018.

\bibitem{altman1999constrained}
E.~Altman.
\newblock \emph{Constrained Markov Decision Processes}.
\newblock Chapman \& Hall/CRC, 1999.

\bibitem{amodei2016concrete}
D.~Amodei, C.~Olah, J.~Steinhardt, P.~Christiano, J.~Schulman, and D.~Mané.
\newblock Concrete problems in AI safety.
\newblock \emph{arXiv:1606.06565}, 2016.

\bibitem{gray1978notes}
J.~Gray.
\newblock Notes on data base operating systems.
\newblock In \emph{Operating Systems. An Advanced Course}, Springer, 1978.

\bibitem{jacobson1988congestion}
V.~Jacobson.
\newblock Congestion avoidance and control.
\newblock \emph{ACM SIGCOMM}, 1988.

\bibitem{saltzer1975protection}
J.~H. Saltzer and M.~D. Schroeder.
\newblock The protection of information in computer systems.
\newblock \emph{Proceedings of the IEEE}, 63(9):1278--1308, 1975.


\bibitem{abortedkv2026}
Guijia Zhang and Harry Yang.
``Aborted but Not Forgotten: KV-Cache Retention Breaks Rollback Consistency in Language Agents.''
arXiv:2608.15939, 2026.

\bibitem{saferesume2026}
Guanlong Wu et al.
``Safe to Resume? Breaking Execution Continuity of Agent Execution via Rollback.''
arXiv:2608.29381, 2026.
\bibitem{garciamolina1987sagas}
H. Garcia-Molina and K. Salem.
\newblock Sagas.
\newblock In {\em ACM SIGMOD}, 1987.

\bibitem{gray1992transaction}
J. Gray and A. Reuter.
\newblock {\em Transaction Processing: Concepts and Techniques}.
\newblock Morgan Kaufmann, 1992.

\bibitem{wiener1948}
N.~Wiener.
\newblock \emph{Cybernetics. Or Control and Communication in the Animal and the Machine}.
\newblock MIT Press, 1948.

\end{thebibliography}
% ============================================================
